\clearpage
\onecolumn
\appendix

\section{Algorithm}\label{sec:algorithm}

\begin{algorithm}[h]
    \caption{Dynamic MCTS-based KGQA with Path Model Pretraining and Online Refinement}
    \renewcommand{\algorithmicrequire}{\textbf{Input:}}
    \renewcommand{\algorithmicensure}{\textbf{Output:}}
    \begin{algorithmic}[1]
        \REQUIRE Question $q$, knowledge graph $\mathcal{G} = (\mathcal{E}, \mathcal{R}, \mathcal{T})$, number of selected relations $k$, MCTS iterations $N$, length of reasoning path $L$
        \ENSURE Answer set $\mathcal{A}$
        
        \STATE \texttt{/ Stage 1: Path Evaluation Model Pre-training /}
        \STATE Construct reasoning path pairs $(q, p^{+}, p^{-})$ from $\mathcal{G}$ 
        \STATE Initialize path evaluation model $S(q, \cdot; \Theta)$
        \FOR{each batch in pretraining data}
            \STATE Update $S(q, \cdot; \Theta)$ by minimizing the Pair-wise Ranking loss in Eq.(4) 
        \ENDFOR

        \STATE \texttt{/ Stage 2: Dynamic MCTS Reasoning /}
        \FOR{$i = 1$ to $N$}
            \STATE \textbf{Selection:}
            Traverse the tree from root to a leaf node by selecting child nodes according to the UCT criterion in Eq.(1)
            \STATE \textbf{Expansion:}
            
            \STATE i. At the selected node, enumerate all candidate relations from current entities and use the LLM-based planner to select the top-$k$ most relevant relations
            
            \STATE ii. Expand a new child node for each selected relation
            \STATE \textbf{Simulation:}
            For each expanded node, perform a rollout by sequentially selecting relations (guided by the path evaluation model) up to $L$ hops or until a correct answer is reached
            \STATE \textbf{Backpropagation:}
            Update the value ($w_i$) and visit ($n_i$) statistics along the traversed path from the leaf node back to the root using the score from the simulation, as per Eq.(6)
            
            \STATE \textbf{Path Evaluation Model Fine-tuning:} Generate the explored pseudo-path pairs $(\hat{p}^{+}, \hat{p}^{-})$ via Eq.(5) and fine-tune  the path evaluation model $S(q, \cdot; \Theta)$ based on Pair-wise Ranking loss in Eq.(4)
        \ENDFOR

        \STATE \texttt{/ Stage 3: Answer Extraction /}
        \STATE Collect entities reached by high-scoring reasoning paths as $\mathcal{A}$
        \STATE \textbf{return} $\mathcal{A}$
    \end{algorithmic}
    \label{algo:dynamic-mcts-final}
\end{algorithm}

% \section{Complexity Analysis}

% The overall time complexity of the proposed \method{} framework is governed by its two primary online stages: the LLM-Guided MCTS Search and the interleaved Path-based Dynamic Refinement. In the MCTS Search phase, executed over $N$ iterations, the LLM-guided relation expansion incurs a total complexity of $\mathcal{O}\left(N \cdot T_{\text{LLM}}(k)\right)$, where $T_{\text{LLM}}(k)$ denotes the inference time of the LLM when provided with up to $k$ candidate relations from the Knowledge Graph (KG). The context-aware path evaluation performed during the simulation step introduces an additional cost of $\mathcal{O}\left(N \cdot k \cdot L^3 \cdot d\right)$, where $L$ is the maximum path length and $d$ is the embedding dimension. In the Path-based Dynamic Refinement stage, the evaluator is fine-tuned for $N_{\text{FT}}$ steps, with a total complexity of $\mathcal{O}\left(N_{\text{FT}} \cdot L^2 \cdot d\right)$. Consequently, the overall time complexity of \method{} is:
% $\mathcal{O}\left(N \cdot\left(T_{L L M}(k)+k \cdot L^3 \cdot d\right)+N_{F T} \cdot L^2 \cdot d\right)$.



\section{ Datasets}\label{sec:dataset}

\subsection{Dataset Description}
\begin{table}[h]
\centering
\label{tab:dataset_statistics}
\begin{tabular}{lccc}
\toprule
\textbf{Datasets} & \textbf{\#Train} & \textbf{\#Valid} & \textbf{\#Test} \\
\midrule
WebQSP     & 2,848   & 250    & 1,639  \\
\midrule
CWQ        & 27,639  & 3,519  & 3,531   \\
% \midrule
% MetaQA & 329,282 & 39,136 & 30,903 \\ 
\bottomrule
\end{tabular}
\caption{Statistics of KGQA datasets.}
\label{tab:datasets}
\end{table}


We conduct extensive experiments on two widely used multi-hop Knowledge Graph Question Answering (KGQA) benchmarks: WebQSP~\cite{talmor2018web} and CWQ~\cite{yih2016value}. The statistics of these two benchmarks can be found in Table \ref{tab:datasets}, and
their details are shown as follows:

\begin{itemize}
    \item The WebQuestionsSP (WebQSP) dataset is a widely adopted benchmark for evaluating single-hop and simple multi-hop KGQA~\cite{yih2016value}. It consists of 4,837 natural language questions annotated with corresponding SPARQL queries over the Freebase knowledge graph. The dataset is partitioned into 2,848 training, 250 validation, and 1,639 test instances.
    \item The ComplexWebQuestions (CWQ) dataset is a challenging benchmark designed for multi-hop KGQA~\cite{talmor2018web}. It comprises 34,689 questions derived from WebQuestionsSP, reformulated to include more complex and compositional queries. Each question typically requires multi-step reasoning over the Freebase knowledge graph, often involving conjunctions, comparatives, or nested logical structures. The dataset is divided into 27,639 training, 3,519 validation, and 3,531 test examples.

    % \item The MetaQA dataset serves as a large-scale benchmark for evaluating multi-hop knowledge graph question answering in the movie domain~\cite{zhang2018variational}. It comprises 399,321 automatically generated natural language questions based on the WikiMovies knowledge graph, partitioned into 329,282 training, 39,136 validation, and 30,903 test examples. The dataset is organised into three subsets according to the number of reasoning steps required to answer each question: 1-hop, 2-hop, and 3-hop. Each instance challenges models to perform compositional reasoning over complex entity and relation chains, providing a rigorous assessment of both single-step and multi-step KGQA capabilities.
\end{itemize}

\subsection{Data Processing}

Following prior work~\cite{shen2025reasoning,long2025enhancing, ma2025deliberation}, we preprocess the datasets by constructing localized subgraphs centered around each question entity to reduce the size of the search space. Specifically, for each question in WebQSP~\cite{yih2016value} and CWQ~\cite{talmor2018web}, we extract a subgraph from the Freebase knowledge graph by including all triples within a predefined number of hops from the topic entity. This approach preserves the essential context required for multi-hop reasoning while significantly improving computational efficiency. 

\section{ Baselines}\label{sec:baseline}

In this part, we introduce the details of the compared baselines as follows:

\begin{itemize}

\item \textbf{Semantic Parsing Methods.} We compare our \method{} with six semantic parsing methods:

\begin{itemize}
    \item \textbf{KV-Mem}: KV-Mem~\cite{miller2016key} introduce a neural architecture that stores facts as key-value pairs and enables question answering by attending over memory slots, directly retrieving relevant information to infer answers.
    \item \textbf{EmbedKGQA}: EmbedKGQA~\cite{saxena2020improving} enhances multi-hop question answering over knowledge graphs by leveraging pretrained knowledge base embeddings, enabling the model to reason over entity and relation representations without explicit path enumeration during answer prediction.
    \item \textbf{QGG}: QGG~\cite{lan2020query} generates query graphs to answer multi-hop complex questions over knowledge bases, formulating question answering as query graph prediction and enabling structured reasoning through graph matching and path ranking mechanisms.
    \item \textbf{NSM}: NSM~\cite{he2021improving} enhances multi-hop KBQA by leveraging intermediate supervision signals, decomposing questions into reasoning steps, and training a neural state machine to sequentially predict relations and entities for accurate path-based reasoning.
    \item \textbf{TransferNet}: TransferNet~\cite{shi2021transfernet} proposes a transparent framework for multi-hop QA over relational graphs by transferring question semantics to relation paths through interpretable path ranking and structured reasoning, enabling effective and explainable answer prediction.
    \item \textbf{KGT5}: KGT5~\cite{saxena2022sequence} formulates knowledge graph completion and question answering as unified sequence-to-sequence tasks, leveraging pre-trained language models to jointly encode input queries and generate answer entities or triples in a flexible and end-to-end manner.
    
\end{itemize}

\item \textbf{Retrieval-Based Methods.} We compare our \method{} with four retrieval-based methods:

\begin{itemize}
    \item \textbf{GraftNet}: GraftNet~\cite{sun2018open} proposes an early fusion framework that jointly encodes knowledge base facts and supporting text by constructing a heterogeneous graph, enabling effective reasoning through graph convolutional networks for open-domain question answering.
    \item \textbf{PullNet}: PullNet~\cite{sun2019pullnet} introduces an iterative retrieval mechanism that expands a query-specific subgraph by pulling relevant facts from both knowledge bases and text, enabling joint reasoning over heterogeneous evidence for open-domain question answering. 
    \item \textbf{SR+NSM}: SR+NSM~\cite{zhang2022subgraph} enhances multi-hop KBQA by first retrieving a question-relevant subgraph and then performing symbolic reasoning over it using Neural Symbolic Machines, improving efficiency and accuracy through constrained and focused logical inference.
    \item \textbf{SR+NSM+E2E}: SR+NSM+E2E~\cite{zhang2022subgraph} extends SR+NSM by enabling end-to-end training that jointly optimizes subgraph retrieval and reasoning. This integration enhances model coherence and allows better alignment between retrieved subgraphs and final answer prediction.
\end{itemize}


    \item  \textbf{General
Large Language Models (LLMs).} We compare our \method{} with six general
LLMs:

\begin{itemize}

    \item \textbf{Flan-T5-xl}: Flan-T5-xl~\cite{chung2024scaling} is an instruction-finetuned variant of the T5 model, trained on a diverse collection of tasks with natural language instructions. By leveraging large-scale instruction tuning, it improves zero-shot and few-shot performance across diverse NLP benchmarks.
    \item \textbf{Alpaca-7B}: Alpaca-7B~\cite{taori2023stanford} is an instruction-following language model fine-tuned from LLaMA-7B using self-instruct techniques. It demonstrates strong zero-shot and few-shot performance by aligning with human instructions across various NLP tasks.
    \item \textbf{Llama3-8B}: Llama3-8B~\cite{dubey2024llama} is part of the LLaMA 3 family of models, designed for improved instruction following, reasoning, and code generation. Pretrained on a high-quality corpus and fine-tuned with supervised signals, it achieves strong performance across diverse benchmarks.
    \item \textbf{Qwen2.5-7B}: Qwen2.5-7B~\cite{team2024qwen2} is a 7B-parameter instruction-tuned language model developed by Alibaba, optimized for tasks such as reasoning, code generation, and dialogue. It supports multi-turn conversation and demonstrates competitive performance on standard benchmarks.
    \item \textbf{ChatGPT}: ChatGPT~\cite{schulman2022chatgpt} is a conversational AI developed by OpenAI, based on the GPT architecture. It is designed to understand natural language, engage in dialogue, answer questions, and assist with a wide range of tasks across domains.
    \item \textbf{ChatGPT+CoT}: ChatGPT with Chain-of-Thought (CoT)~\cite{wei2022chain} prompting enhances the model's reasoning capabilities by encouraging it to generate intermediate reasoning steps before arriving at a final answer, improving performance on complex, multi-step problems.
\end{itemize}


    \item \textbf{LLMs with KG.} We compare our \method{} with fourteen LLMs with KG methods:

\begin{itemize}
    \item  \textbf{UniKGQA}: UniKGQA~\cite{jiang2022unikgqa} is a unified framework that integrates retrieval and reasoning for multi-hop question answering over knowledge graphs, combining subgraph retrieval, query decomposition, and neural reasoning in an end-to-end manner.
    \item \textbf{DECAF}: DECAF~\cite{yu2022decaf} is a joint framework for question answering over knowledge bases that simultaneously decodes logical forms and answers. By leveraging dual supervision, it enhances both symbolic reasoning accuracy and direct answer prediction in a unified architecture.
    \item \textbf{KD-CoT}: KD-CoT~\cite{wang2023knowledge} is a framework that enhances the faithfulness of large language models by guiding Chain-of-Thought reasoning with external knowledge, improving accuracy in knowledge-intensive question answering tasks.
    \item \textbf{Nutrea}: Nutrea~\cite{choi2023nutrea} proposes a neural tree search framework for context-guided multi-hop KGQA. It incrementally constructs reasoning trees by integrating question semantics and graph context, enabling efficient exploration of multi-hop paths for accurate answer prediction.
    \item \textbf{ToG}: ToG~\cite{sun2023think} is a framework that enables large language models to perform deep and responsible reasoning over knowledge graphs by combining structured graph information with iterative thinking and verification mechanisms for reliable multi-hop QA.
    \item \textbf{RoG}: RoG~\cite{luo2023reasoning} is a framework that enhances the faithfulness and interpretability of large language model reasoning by grounding multi-hop question answering on knowledge graphs, integrating symbolic path tracking with natural language generation.
    \item \textbf{KAPING}: KAPING~\cite{baek2023knowledge} introduces knowledge-augmented prompting by integrating structured triples into Chain-of-Thought (CoT) reasoning. It guides large language models to generate intermediate reasoning steps, enabling zero-shot multi-hop KGQA without task-specific fine-tuning.
    \item \textbf{ReasoningLM}: ReasoningLM~\cite{jiang2023reasoninglm} enhances pre-trained language models for KGQA by injecting subgraph structures into the input representation. It enables structural reasoning over retrieved subgraphs through a reasoning-aware encoder, improving performance on complex multi-hop queries.
    \item \textbf{FiDeLis}: FiDeLis~\cite{sui2024fidelis} proposes a faithfulness-aware KGQA framework that enhances reasoning consistency in LLMs by aligning generated logical forms with answer predictions. It introduces fidelity constraints to reduce hallucinations and improve answer correctness.
    
    % \item \textbf{PoG}: PoG~\cite{chen2024plan} is a self-correcting framework that enables large language models to perform adaptive planning over knowledge graphs by dynamically revising reasoning steps, enhancing robustness and accuracy in multi-hop question answering.
    
    % \item \textbf{RoG}: G-Retriever~\cite{he2024g} is a retrieval-augmented generation framework designed for textual graph understanding and question answering, which integrates graph-structured retrieval with language modeling to enhance context relevance and reasoning accuracy.
    %\item \textbf{EtD}: EtD~\cite{liu2024explore} is a GNN-LLM synergy framework for knowledge graph reasoning, where a GNN explores plausible reasoning paths and an LLM determines the final answer, combining structured exploration with semantic understanding.
    % \item \textbf{GNN-RAG}: GNN-RAG~\cite{mavromatis2024gnn} is a framework that integrates Graph Neural Networks with retrieval-augmented generation, enabling large language models to reason more effectively by retrieving and encoding structured knowledge from graphs for informed answer generation.
    \item \textbf{GNN-RAG}: GNN-RAG~\cite{mavromatis2024gnn} integrates graph neural networks with retrieval-augmented generation by encoding knowledge subgraphs into LLMs' context. It enables structural reasoning over retrieved subgraphs, improving answer accuracy in KGQA through explicit graph-aware representations.
    %\item \textbf{Interactive-KBQA}: Interactive-KBQA~\cite{xiong2024interactive}  is a framework that enables multi-turn interactions between large language models and knowledge bases, allowing iterative question refinement and retrieval to improve accuracy and coherence in complex knowledge-based question answering.
    \item \textbf{DoG}: DoG~\cite{ma2025debate} is a flexible and reliable reasoning framework that enables large language models to generate and evaluate multiple reasoning paths over knowledge graphs through a debate-style process, enhancing robustness and answer faithfulness.
    %\item \textbf{SubgraphRAG}: SubgraphRAG~\cite{ma2025knowledge} is a framework that leverages knowledge graph retrieval to provide structured context for large language models, enhancing their ability to align and match heterogeneous schemas accurately and interpretably.
    \item \textbf{DuarL}: DuarL~\cite{liu2025dual} is a collaborative framework that integrates GNNs and LLMs for KGQA, where GNNs capture structural semantics and LLMs perform adaptive reasoning, enabling accurate and interpretable multi-hop QA.
    \item \textbf{DP}: DP~\cite{ma2025deliberation} is a trustworthy reasoning framework that guides large language models using prior knowledge from knowledge graphs. It iteratively verifies and refines reasoning paths to enhance faithfulness, robustness, and answer accuracy in KGQA.
    %\item \textbf{RAPL}: RAPL~\cite{yao2025learning} is a framework designed for KGQA, which learns a lightweight and transferable graph retriever to efficiently select task-relevant subgraphs, enhancing both scalability and generalization across diverse questions.
    \item \textbf{RwT}: RwT~\cite{shen2025reasoning} is a faithful KGQA framework that models multi-hop reasoning as tree-structured exploration over knowledge graphs, enabling large language models to generate interpretable reasoning paths and improve answer consistency and accuracy.

\end{itemize}

\end{itemize}

\section{More experimental results}\label{sec:experiments_results}

To more thoroughly illustrate the impact of hyperparameter variations on model performance, we report detailed numerical results showing how performance fluctuates under different hyperparameter settings. As presented in Table~\ref{tab:sensitive_k} and Table~\ref{tab:sensitive_length}, these results provide a comprehensive understanding of the model’s sensitivity and stability across a range of configurations.

\begin{table}[t]
\centering
\small
\begin{tabular}{lcccc}
\toprule
\multirow{2}{*}{Method} & \multicolumn{2}{c}{WebQSP} & \multicolumn{2}{c}{CWQ} \\
 & Hits@1 & F1 & Hits@1 & F1 \\
\midrule
$k=2$ & 93.0 & 80.9 & 76.6 & 73.8 \\
$k=3$ & 94.0 & 81.7 & 78.0 & 75.1 \\
$k=4$ & 94.0 & 81.8 & 77.8 & 75.0 \\
$k=5$ & 93.9 & 81.7 & 78.0 & 75.2 \\
\bottomrule
\end{tabular}
\caption{Hyperparameter sensitivity analysis of the number of selected relations $k$ on the WebQSP and CWQ datasets.}
\label{tab:sensitive_k}
\end{table}



\begin{table}[t]
\centering
\small
\begin{tabular}{lcccc}
\toprule
\multirow{2}{*}{Method} & \multicolumn{2}{c}{WebQSP} & \multicolumn{2}{c}{CWQ} \\
 & Hits@1 & F1 & Hits@1 & F1 \\
\midrule
$L=2$ & 93.6 & 81.2 & 77.4 & 74.5 \\
$L=3$ & 94.0 & 81.7 & 77.6 & 74.7 \\
$L=4$ & 93.7 & 81.4 & 78.0 & 75.1 \\
$L=5$ & 93.8 & 81.6 & 77.9 & 74.9 \\
\bottomrule
\end{tabular}
\caption{Hyperparameter sensitivity analysis of the reasoning path length $L$ on the WebQSP and CWQ datasets.}
\label{tab:sensitive_length}
\end{table}



% \begin{table*}[t]
% \renewcommand{\arraystretch}{1.3}
% \small
% \centering
% % \vspace{-0.1cm}
% \begin{tabular}{>{\centering\arraybackslash}m{2.5cm}|p{7.6cm}}
% \toprule
% \textbf{Question} & Where does greek language come from? \\
% \hline
% \textbf{Answer} & \textcolor{red}{Attic Group} \\
% \hline
% \multirow{5}{*}{\textbf{Llama-2-13B}} & The Greek language originates from \textcolor{blue}{the Proto-Indo-European language family}, which spread to Greece around 2000 BCE. It evolved into the ancient Greek language, influenced by Mycenaean and Aeolic dialects, and eventually developed into Modern Greek. \\
% \hline
% \multirow{5}{*}{\textbf{Qwen-3-14B}} & The Greek language originates from \textcolor{blue}{the Proto-Hellenic branch of the Indo-European language family}. It evolved from the Mycenaean Greek of the second millennium BCE and became the foundation for Classical Greek, spoken in ancient Greece. \\
% \hline
% \multirow{5}{*}{\textbf{GPT 4.1}} & The Greek language originates from \textcolor{blue}{the Indo-European language family}. It developed in the southern Balkans around the 3rd millennium BCE and is the oldest recorded living language, with written records dating back to the 15th century BCE (Linear B script). \\
% \hline
% \multirow{5}{*}{\textbf{GPT 4.1-mini}} & The Greek language originates from \textcolor{blue}{the Hellenic branch of the Indo-European language family}. It developed in the region of Greece and the Aegean around the 3rd millennium BCE, evolving through ancient, medieval, and modern stages. \\
% \hline
% \multirow{2}{*}{\textbf{\method{}}} & \texttt{Path 1}: Entity (id: 120026) $\rightarrow$ base.rosetta.languoid.parent $\rightarrow$ \textcolor{red}{Attic Group}. \\
% \bottomrule
% \end{tabular}
% \caption{Case study of \method{}. We highlight the correct answers in \textcolor{red}{Red} and the wrong answers in \textcolor{blue}{Blue}.}
% \label{tab:case_study_method_2}
% \end{table*}


% \subsection{More case study}\label{sec:case_study}

% Table~\ref{tab:case_study_method_2} presents a case study comparing the answer accuracy of \method{} with general LLMs: Llama-2-13B, Qwen-3-14B, GPT 4.1-mini, and GPT 4.1. When asked about the origin of the Greek language, all baseline models generate fluent and seemingly plausible responses grounded in general linguistic knowledge, such as “Proto-Indo-European” or “Proto-Hellenic”, but fail to identify the correct answer: Attic Group. In contrast, \method{} accurately predicts the correct entity by explicitly traversing the relation path \textit{base.rosetta.languoid.parent} within the knowledge graph. This example illustrates a key advantage of \method{}: rather than relying solely on learned linguistic patterns, it performs structured reasoning over the knowledge graph, enabling precise and faithful answers to ontology-specific queries that often elude general-purpose LLMs.



\section{Prompt Template}

We provide the prompt templates used by the LLM-based planner to select the top-$k$ most relevant relations from the candidate set at each step of path expansion in Fig. 3, as part of the LLM Guided Path Expansion module.

\begin{tcolorbox}[title=Prompt Template for LLM-Guided Path Expansion, colback=white, colframe=black!75, boxrule=0.5pt, arc=2mm]

\textbf{Role} \\
You are an expert assistant for Knowledge Graph Question Answering (KGQA). Your core capability is to deeply understand natural language questions and the semantics of knowledge graph relations to find the most relevant reasoning paths.

\textbf{Task} \\
Your task is to act as a \textbf{“Relation Retriever.”} Given a natural language question and a list of candidate relations, you must analyze the semantics of the question and each relation to select up to \texttt{k} relations that are most likely to lead to the correct answer.


\textbf{Rules and Constraints}
\begin{itemize}[leftmargin=1.5em, itemsep=0pt, topsep=0pt]
    \item \textbf{Fidelity to Candidates}: Your selection of relations \textbf{MUST} come strictly from the provided \texttt{Candidate Relations} list. Do not invent or modify relations.
    \item \textbf{Quantity Limit}: Return no more than \texttt{k} relations. If multiple relations are highly relevant, order them from most to least relevant. If there are fewer than \texttt{k} relevant relations, return only those.
    \item \textbf{Output Format}: Your response \textbf{MUST} be a list of strings, containing the names of the relations you have selected.
\end{itemize}

\textbf{Example}
\begin{itemize}[leftmargin=1.5em, itemsep=0pt, topsep=2pt]
    \item \textbf{Input:}
    \begin{itemize}
        \item \texttt{Question}: "who was the president after jfk died"
        \item \texttt{Candidate Relations}: \{"government.president", "government.president.successor", "location.location.containedby", "people.person.place\_of\_birth"\}
        \item \texttt{K}: 2
    \end{itemize}
    
    \item \textbf{Output:}
\begin{verbatim}
["government.president", "government.president.successor"]
\end{verbatim}
\end{itemize}

\textbf{Your Task}
\begin{itemize}[leftmargin=1.5em, itemsep=0pt, topsep=2pt]
    \item \texttt{Question}: \{question\}
    \item \texttt{Candidate Relations}: \{relations\_list\}
    \item \texttt{K}: \{k\}
\end{itemize}

\textbf{Output:}
\begin{verbatim}
[]
\end{verbatim}
\end{tcolorbox}
\captionof{figure}{Prompt template used in the LLM-based planner to select top-k relations during reasoning.}
\label{fig:prompt_llm_planner}


\section{The Use of Large Language Models (LLMs)}

Large language models (LLMs) were only used to improve the clarity, grammar, and fluency of the manuscript. They were not involved in the development of research ideas, experimental design, data analysis, or any other aspect of the scientific content.